% The spoken body of the homily for the resumed Fifth Sunday after the
% Epiphany, said on Sunday 8 November 2026. Continuous preaching, ready to read
% aloud after the Gospel: no heading, no direction to the preacher, no citation
% inside the speech. The loci, the relation to the reviewed interpretations,
% the audience, the word count and the pace estimate stand in the terminal
% note, which no one reads out. The blank lines are the paragraph breaks a
% speaker recovers his place by. The speech tells the parable; gives the Lord's
% own explanation, with the good seed as the children of the kingdom; sets the
% servants' offer beside St Paul's command to the Corinthians, as St Jerome
% does; gives the two things the Fathers find in the householder's reason, the
% likeness of the plants and the room for change; keeps the Church's
% correction; gives the servants' work from the Epistle, St Rabanus and the
% letter from which the Introit comes; turns to the would-be weeders through
% the Alleluia and Offertory, the Secret and St Augustine's wish; brings the
% Collect and Gradual in with St Thomas; gives one response from the
% Communion's next verse and St Bede; and closes on the harvest and the
% Postcommunion. The honorific is tied to the name so that ``St'' never ends a
% line.
\raggedcolumns

The Gospel we have just heard is a parable about a field. A man sows good
seed in it. While men are asleep, his enemy comes, sows cockle among the
wheat, and goes his way. For a time nothing shows. Only when the blade has
sprung up and borne fruit does the cockle appear. Then the servants come to
the householder: ``Sir, didst thou not sow good seed in thy field? Whence then
hath it cockle?'' He answers, ``An enemy hath done this.'' They ask whether
they should go and gather it up, and he tells them no.

Afterwards, in the house, the Lord explained the parable to his disciples. The
field, he said, is the world, and ``the good seed are the children of the
kingdom''; the cockle are the children of the wicked one. The enemy is the
devil, the harvest is the end of the world, and the reapers are the angels.
So the wheat and the cockle are people. St~Augustine, preaching on this
Gospel, took the cockle for bad Christians among the good. Other translations
call it tares, and he would not let his people think the tares grew only
among them. ``Even in these high seats,'' he told them, ``there is both wheat,
and tares, and among the laity there is wheat, and tares.'' The high seats
were the bench where he sat as bishop with his priests.

The servants' offer is a sensible one, and St~Paul seems to stand with them.
He told the Corinthians to expel a man living in incest. St~Jerome set Paul's
command, ``Take away the evil from among you,'' beside the householder's
``Suffer both to grow,'' and asked how the Church can be told both.

The householder gives his reason in the parable itself: ``lest perhaps
gathering up the cockle, you root up the wheat also together with it.'' The
Fathers find two things in that reason.

The first is the plant. The cockle here is the weed St~Jerome calls darnel,
and he says that while darnel and wheat are still in the blade, the two can
hardly be told apart, or not at all. The parable says the same: the cockle
appeared only when the wheat bore fruit. A man who went into the field to
weed while the crop was green would be pulling by guesswork. St~Augustine
counts himself and his hearers among the servants, and gives the reason they
were refused: ``ye are such as may be deceived.'' We are men, he says, and the
reapers are the angels. Jerome draws his answer to his own question from the
same likeness. Where a case is doubtful, we are not to pass sentence quickly,
but to leave the judgment to God, so that what God removes at the end is
manifest guilt and not a suspicion.

The second thing is in the words ``lest perhaps.'' No farmer expects darnel to
turn into wheat, but the cockle of this parable is men. Jerome reads in
``lest perhaps'' that room is given for repentance, because one who today is
corrupted by false teaching may tomorrow come to his senses; St~John
Chrysostom reads the same there. Augustine, in the same sermon, turned to the
tares in his own church and told them to amend before the harvest came. So
the householder waits for the cockle's sake as well as the wheat's: he leaves
it time to change.

The Church still corrects. Augustine wrote that until the end of the world she
goes on ``correcting as many as she can, and bearing with those whom she
cannot correct.'' The separation itself waits for the harvest.

The householder gives his servants no task but to let both grow. For what
they are to do meanwhile, take St~Paul's words in today's Epistle: ``bearing
with one another and forgiving one another, if any have a complaint against
another. Even as the Lord hath forgiven you, so do you also.'' St~John
Chrysostom says that Paul's word ``complaint'' makes the offence a small
thing; but ``even as the Lord hath forgiven you'' means that we are to forgive
even serious charges, and not only to forgive but to love.

St~Rabanus Maurus finds the same patience in the householder himself. He had
been wronged, and he knew who had done it: ``An enemy hath done this.'' Yet he
did not blaze up in anger or avenge himself at once; he waited. So, Rabanus
says, when we have received an injury from our enemies, we ought not to break
out at once into revenge, but to have recourse to patience, and to keep it to
the end of life.

There is one more thing the servants can do, and it comes from the letter
that gave us today's Introit. The prophet Jeremias wrote that letter from
Jerusalem to the Jews carried off to Babylon. Before the promise the Introit
sings, ``I think thoughts of peace, and not of affliction,'' the letter tells
the exiles to ``seek the peace of the city, to which I have caused you to be
carried away captives; and pray to the Lord for it.'' St~Jerome reads that
verse of the Church: we are commanded to pray for our enemies, ``so that we may not be content with our
own salvation alone, but seek our enemies' salvation as well.'' That prayer
can take in the cockle the servants are not to pull up, and ask for the change
Augustine asked of it.

Before the Gospel, the Alleluia verse was the opening of the hundred and
twenty-ninth psalm: ``From the depths I have cried to Thee, O Lord: Lord, hear
my prayer.'' At the Offertory it will be sung again. St~Augustine asks of that
psalm, ``Who is it then who crieth?'' and he answers, ``A sinner.'' We sing
that cry ourselves. By his answer, those of us who would like the field
cleared are sinners asking pardon; in the Secret the priest will ask God to
``mercifully forgive our sins.'' The wish that Augustine puts among the Lord's
words to the zealous servants is therefore meant for us as well: ``That time
will come, may it only find you wheat!''

St~John Chrysostom says the servants' haste showed their anxiety for the
crop. In the Collect we brought that anxiety to God, asking him to keep his
family with continual mercy, so that it may always be defended by his
protection. The Gradual thanked him: ``Thou hast delivered us, O Lord, from
them that afflict us.'' St~Thomas Aquinas, expounding that psalm, says that
the saints are afflicted in two ways, in the body, and in the spirit when
those who live among them sin; and from these afflictions, he says, God frees
them. That second affliction is the wheat's, growing beside the cockle. It is
God, and not the servants' hands, who frees the wheat from the trouble the
cockle gives it.

The Communion antiphon is the Lord's promise in St~Mark: ``whatsoever you ask
when you pray, believe that you shall receive, and it shall be done unto
you.'' The next verse in Mark, which the antiphon does not sing, goes on:
``And when you shall stand to pray, forgive, if you have aught against any
man: that your Father also, who is in heaven, may forgive you your sins.''
St~Bede reads the two verses together. Those who are not yet perfect, he says,
will obtain what they ask, ``provided that they first forgive others who sin
against them.'' So before you come to the altar today, forgive the one
against whom you have a complaint, as the Lord has forgiven you, and pray for
him.

The harvest will come. To the good in his church, who grieved to be few among
many, Augustine promised that ``the angels will come who can make the
separation, and who cannot make mistakes.'' The Lord's own explanation of
the parable ends, ``Then shall the just shine as the sun, in the kingdom of
their Father.'' At this altar you will
receive what the Postcommunion calls the pledge of salvation, and the priest
will then ask that we may receive its effect. The pledge is given to us now,
in the field, while the wheat and the cockle grow together. Its effect is
what the householder promised the wheat: ``gather ye into my barn.''
